\section{Preliminaries}
\label{sec:a-pre-preliminaries}

The algorithms below separate two questions: how the graph restricts an
optimizing scenario, and how accurately its physical state can be
computed from binary input. This section fixes the passive-flow model,
proves the block decomposition used in the structural arguments, and
states the arithmetic and output conventions that separate exact
comparison from high-precision approximation.

\subsection{Networks, laws, and passive states}
\label{subsec:a-pre-model}
Let $G=(V,E)$ be a finite connected loopless graph, with $n=|V|$ and
$m=|E|$. Parallel edges are allowed unless simplicity is stated.
Choose an orientation for each edge. An arc $e=(u,v)$ has tail $u$ and
head $v$; its column in the incidence matrix $A\in\R^{V\times E}$ has
$+1$ at $u$ and $-1$ at $v$. A nomination $b\in\R^V$ specifies injections
(positive entries) and withdrawals (negative entries), and is balanced:
$\one^\top b=0$. Flows $x\in\R^E$ are signed relative to the orientation.
A passive state consists of a flow and potentials $\pi\in\R^V$ satisfying
\begin{equation}
 Ax=b,\qquad A^\top\pi=g(x):=(g_e(x_e))_{e\in E}.
 \label{eq:a-pre-passive}
\end{equation}
A general edge law $g_e:\R\to\R$ is continuous and strictly increasing,
with $g_e(0)=0$; neither surjectivity nor a growth condition is assumed. The symmetric quadratic law is $g_e(t)=\beta_e t|t|$, with
$\beta_e>0$. The asymmetric quadratic law has two positive coefficients:
\[
 g_e(t)=\begin{cases}\beta_e^+t^2,&t\ge0,\\-\beta_e^-t^2,&t\le0.\end{cases}
\]
The two formulas agree at zero. Reversing an edge replaces $x_e$ by
$-x_e$ and its law by $\widetilde g_e(t)=-g_e(-t)$, so symmetric
coefficients are unchanged and asymmetric coefficients are swapped.

\begin{theorem}[Existence and uniqueness]
\label{thm:a-pre-existence}
For every balanced nomination on a connected graph and every family of
general edge laws as above, a passive state exists. Its flow is unique,
and its potentials are unique up to an additive constant. The flow is
exactly the minimizer of the primitive energy
\begin{equation}
 \energy(x)=\sum_{e\in E}F_e(x_e),\qquad
 F_e(t)=\int_0^t g_e(u)\,du,
 \quad\text{over }\{x:Ax=b\}.
 \label{eq:a-pre-energy}
\end{equation}
\end{theorem}
\begin{proof}
Choose a rooted spanning tree. Route each subtree's total nomination
through its parent edge, with the sign prescribed by the orientation,
and set other flows to zero. Conservation holds at every nonroot vertex;
balance gives it at the root. This constructs a feasible $x^0$.

Each $F_e$ is continuously differentiable with strictly increasing
derivative $g_e$, hence strictly convex. Also $F_e(t)\ge0$. For $t\ge1$,
\[
 F_e(t)\ge g_e(1)(t-1),
\]
and for $t\le-1$,
\[
 F_e(t)=-\int_t^0g_e(u)\,du
       \ge[-g_e(-1)](|t|-1).
\]
With $a_e=\min\{g_e(1),-g_e(-1)\}>0$ we get
$F_e(t)\ge a_e|t|-a_e$ for all $t$, including $|t|\le1$. If $m>0$,
put $a=\min_e a_e>0$; then
\[
 \energy(x)\ge a\norm{x}_1-\sum_e a_e.
\]
For $x=x^0+z$ with $z\in\ker A$, this lower bound is at least
$a\norm{z}_1-a\norm{x^0}_1-\sum_e a_e$, so sublevel sets on the closed
affine feasible space are bounded. The nonempty sublevel set at
$\energy(x^0)$ is compact, and continuity gives a minimizer. Strict convexity gives uniqueness. If $m=0$, the
connected graph has one vertex and the assertion is immediate.

At the minimizer every directional derivative along $\ker A$ vanishes,
so $g(x)\in(\ker A)^\perp=\operatorname{im}A^\top$, giving a
multiplier $\pi$ with $A^\top\pi=g(x)$. No constraint qualification is
needed, since $(\ker A)^\perp=\operatorname{im}A^\top$ holds for every
matrix, including the rank-deficient incidence matrix.
Conversely, if \cref{eq:a-pre-passive} holds, convexity gives, for every
feasible $y$,
\[
 \energy(y)\ge\energy(x)+g(x)^\top(y-x)
 =\energy(x)+\pi^\top A(y-x)=\energy(x).
\]
Finally, $A^\top(\pi-\pi')=0$ makes $\pi-\pi'$ equal at adjacent
vertices, hence constant by connectivity.
\end{proof}

The variational principle is classical. In resistance variables,
\citet[Section 5, Theorems $2'$ and $3'$]{birkhoff1956-non-linear-network-problems}
characterize the Neumann solution by stationarity of the sum of
primitives and obtain existence when those primitives tend to $+\infty$
in both directions; the bound above verifies that condition here.
\citet[Theorems 2 and 3]{duffin1947-nonlinear-networks-iia} proves existence
and uniqueness in conductivity variables (current as a function of
potential drop). His hypotheses apply when each $g_e$ is onto $\R$, for then the
conductivity $g_e^{-1}$ is defined on all of $\R$, continuous, strictly
increasing, and unbounded in both directions; we do not assume this.
\citet[Theorem 6.4 with its Corollary I, and the corollary to Theorem 8.1]{minty1960-monotone-networks}
works more generally with maximal monotone resistor relations.
By Corollary I, any two solutions agree in current and voltage on each
``real resistor'' branch, whose relation contains no horizontal or
vertical segment. His existence criterion is the absence of unbalanced
cycles and cocycles in the associated admissible voltage and current
intervals. The proof above is self-contained for the present
prescribed-nomination model.

\begin{remark}[Bounded laws and normalization]
\label{rem:a-pre-gauge}
Bounded laws do not obstruct existence even with a nontrivial cycle
space. Take two parallel edges from $s$ to $t$ with laws
$g_1(u)=\arctan u$ and $g_2(u)=2\arctan u$, whose suprema are different.
For nomination $(B,-B)$, conservation gives flows $q$ and $B-q$.
Their drops agree exactly when
\[
 \arctan q-2\arctan(B-q)=0.
\]
For every real $B$ the left side is continuous and strictly increasing
in $q$. As $q\to+\infty$, both $\arctan q$ and $-\arctan(B-q)$
have positive limits; as $q\to-\infty$, both have negative limits.
The left side therefore has exactly one root. This is the general-law cycle
equation of \cref{lem:a-pre-cycle} after reversing the second edge.
Superlinear primitive growth, such as the cubic primitive of a quadratic
law, also gives coercivity directly. Potentials are gauge invariant:
adding $\alpha\one$ changes no equation. We fix a reference vertex $v_0$
and impose $\pi_{v_0}=0$ whenever a potential vector is output.
\end{remark}

At every passive state the dissipation identity is
\begin{equation}
 b^\top\pi=(Ax)^\top\pi=x^\top A^\top\pi
           =\sum_e x_e g_e(x_e).
 \label{eq:a-pre-dissipation}
\end{equation}
For symmetric quadratic laws,
\[
 b^\top\pi=\sum_e\beta_e|x_e|^3,\qquad
 \energy(x)=\frac13\sum_e\beta_e|x_e|^3.
\]
For asymmetric quadratic laws use
$\beta_e^+$ on positive flows and $\beta_e^-$ on negative flows in both
sums. In either quadratic model, dissipation is three times the
primitive energy.

\begin{lemma}[Homogeneity]
\label{lem:a-pre-homogeneity}
Fix symmetric or asymmetric quadratic laws. If $(x,\pi)$ is the state
for $b$, then for
$q\ge0$ the nomination $qb$ has flow $qx$ and potential differences
$q^2$ times those for $b$.
\end{lemma}
\begin{proof}
For $q>0$ scaling preserves the sign of each nonzero flow, and both
law pieces vanish at $q=0$. Thus in either quadratic model
$A(qx)=qb$ and $g_e(qx_e)=q^2g_e(x_e)$, so $(qx,q^2\pi)$ is a state.
Uniqueness in \cref{thm:a-pre-existence} gives the claim, also for $q=0$.
\end{proof}

\subsection{Structure lemmas}
\label{subsec:a-pre-structure}
A block is a maximal biconnected subgraph, with every bridge included as
a two-vertex block. Blocks have edges; the one-vertex graph has no blocks.
The block--vertex incidence graph has a node for each block and each
original vertex, with an edge whenever a vertex belongs to a block.
For $n>1$, deleting the nodes of noncut vertices (these are leaves by
\cref{lem:a-pre-incidence}) gives the block--cut tree.

\begin{lemma}[Block incidence tree]
\label{lem:a-pre-incidence}
Every edge belongs to exactly one block, every simple cycle lies in one
block, and distinct blocks share at most one vertex. The block--vertex
incidence graph of a connected graph is a tree, so it defines a unique
block path between any two vertices.
\end{lemma}
\begin{proof}
A cycle is biconnected and extends to a maximal biconnected subgraph
because the graph is finite. Every edge on no cycle is a bridge:
otherwise an alternative path between its endpoints would close a
cycle. Thus every edge belongs to a block. Two distinct blocks cannot
share two vertices: otherwise deleting any one vertex from their union
would leave each block connected (a bridge block leaves a connected
singleton) and would leave a shared vertex, so the union would be
biconnected, contrary to maximality. This also excludes shared edges,
so each edge has a unique block.

Connectivity of $G$ gives connectivity of its incidence graph: replace
each edge along a graph path by the two incidences with its block.
Suppose the incidence graph has a simple cycle with distinct block
nodes $H_1,\ldots,H_j$ and distinct vertex nodes
$w_1,\ldots,w_j$, where $H_i$ contains $w_{i-1},w_i$ cyclically.
The preceding intersection property gives $j\ge3$. The union of these
blocks remains connected after deletion of any vertex $w$: each
$H_i-w$ is nonempty and connected, and at most one of the distinct
joining vertices $w_i$ has been deleted. The remaining joins connect
all these subgraphs in a path or a cycle. Their union is therefore
biconnected, contradicting the maximality of the distinct $H_i$.
This proves acyclicity; if $n=1$, the incidence graph is a single vertex.
Finally, let $n>1$ and suppose $v$ belongs to just one block $H$.
Every edge incident to $v$ lies in $H$, since an edge's block contains
both endpoints and $H$ is the only block containing $v$. Thus any path
between other vertices that passes through $v$ enters and leaves through
neighbors in $H-v$. Rerouting that segment inside the connected graph
$H-v$ gives a walk avoiding $v$, and hence a path avoiding $v$.
Therefore $v$ is not a cut vertex. A vertex in several blocks is a cut
vertex, since a path between the incident branches after its deletion
would contradict the incidence tree. Thus the noncut vertices are
precisely the original-vertex leaves.
\end{proof}

\begin{lemma}[Block decomposition]
\label{lem:a-pre-block}
For the general laws of \cref{thm:a-pre-existence} the following hold.
\begin{enumerate}
\item If $e=(u,v)$ is a bridge and $S$ is the component of $G-e$
containing its tail $u$, then $x_e=\sum_{w\in S}b_w$.
\item Fix a block $H$. Delete its edges, retaining all vertices, and let
$S_{H,v}$ be the component of $G-E(H)$ containing $v\in V(H)$.
These components contain exactly one vertex of $H$ each and partition
$V$. Define the aggregated nomination by
\begin{equation}
 b^H_v=\sum_{w\in S_{H,v}}b_w
      =b_v+\sum_{w\in S_{H,v}\setminus\{v\}}b_w.
 \label{eq:a-pre-aggregation}
\end{equation}
Then $x|_{E(H)}$ is the unique passive flow on $H$ for $b^H$, and
$\pi|_{V(H)}$ is a potential vector for that block.
\item Let the block path from $s$ to $t$ traverse $H_1,\ldots,H_j$
with successive boundary vertices $v_0=s,v_1,\ldots,v_j=t$.
If $\pi^{H_i}$ is any potential vector for $H_i$ at $b^{H_i}$, then
\begin{equation}
 \pi_s-\pi_t=\sum_{i=1}^j
       (\pi^{H_i}_{v_{i-1}}-\pi^{H_i}_{v_i}).
 \label{eq:a-pre-block-path}
\end{equation}
For $s=t$ the sum is empty.
\end{enumerate}
\end{lemma}
\begin{proof}
Summing conservation over the tail component of a bridge cancels every
internal edge and leaves $x_e$, proving (a).
If a component of $G-E(H)$ contained two vertices of $H$, choose a
path in it and take the segment between two consecutive visits to $H$.
This segment has no internal vertex in $H$ and uses no edge of $H$.
Closing it by a path inside $H$ gives a simple cycle with an edge
outside $H$ and an edge inside $H$. By \cref{lem:a-pre-incidence},
every cycle lies in one block and every edge has a unique block, so
the outside edge would belong to $H$, a contradiction. For a bridge
$H$, the same cycle would contain the bridge, which is impossible.
Connectivity ensures that every component meets $H$. Thus the sets
in (b) partition $V$, and $b^H$ is balanced. The aggregate $b^H_v$
includes all branches hanging at $v$ away from $H$ and counts $b_v$
exactly once, even when $v$ belongs to several blocks.

Solve each block with its balanced $b^H$, using
\cref{thm:a-pre-existence}, and assemble their edge flows. We check
global conservation without assuming the restriction property. At a
noncut vertex $v$, its only block has $b^H_v=b_v$. At a cut vertex $v$,
let $H_1,\ldots,H_\varkappa$ be its incident blocks. The components $U_i$ of
$G-v$ are indexed by these blocks; write $B_i=\sum_{w\in U_i}b_w$.
Balance gives $b_v+\sum_i B_i=0$, while
\[
 b^{H_i}_v=b_v+\sum_{h\ne i}B_h=-B_i.
\]
Hence the assembled divergence at $v$ is
$\sum_i b^{H_i}_v=-\sum_i B_i=b_v$.
Root the block--cut tree at any block. Shift the potential constant on
each successive block to match the already assigned potential at its
parent cut vertex; since the incidence graph is a tree, no conflict
arises. All edge laws now hold globally. Global uniqueness identifies
this constructed state with the given state and proves (b), including
the potential restriction claim.
Finally, the shifted block differences telescope along the unique
block path, proving (c). The one-vertex case is immediate.
\end{proof}

\begin{lemma}[Cactus cycle equation]
\label{lem:a-pre-cycle}
Let $C$ be a cycle block of a cactus with general laws as in
\cref{thm:a-pre-existence}, and use its aggregated nomination $b^C$. Orient the cycle consistently
and designate one cycle edge $e_0$ as the chord relative to the spanning
path $C-e_0$. There are linear functions $\lambda_e$ of $b^C$, with
$\lambda_{e_0}=0$, such that every conservation-feasible flow is
$x=\lambda+q\one$. A potential exists for this flow if and only if
\[
 \Phi(q):=\sum_{e\in E(C)}g_e(q+\lambda_e)=0.
\]
The continuous function $\Phi$ is strictly increasing and has exactly
one real root. For symmetric quadratic laws the specialization is
\[
 \Phi(q)=\sum_{e\in E(C)}\beta_e(q+\lambda_e)|q+\lambda_e|.
\]
\end{lemma}
\begin{proof}
Route $b^C$ on the spanning path with zero chord flow. Subtree sums show
that its flows $\lambda_e$ are linear in $b^C$. The kernel of the consistently
oriented cycle incidence matrix is the span of the all-ones edge vector:
zero divergence equates the two incident circulation values at each
vertex. Thus every feasible flow is $\lambda+q\one$. A potential exists
exactly when the sum of the oriented drops around the cycle is zero:
necessity follows by telescoping, and sufficiency by integrating drops
along the spanning path and checking the chord. This condition is
$\Phi(q)=0$.
Every summand is continuous and strictly increasing. In the extended
real sense,
\[
 \lim_{q\to+\infty}\Phi(q)=\sum_e\sup_{u\in\R}g_e(u)>0,\qquad
 \lim_{q\to-\infty}\Phi(q)=\sum_e\inf_{u\in\R}g_e(u)<0.
\]
Indeed $\sup g_e\ge g_e(1)>0$ and $\inf g_e\le g_e(-1)<0$ for each
edge. All terms in each limit have the same strict sign, so the finite
sums are well defined even when infinite. Hence a finite sign-changing
bracket exists, also when all laws are bounded, and the intermediate
value theorem and strict increase prove the assertion.
\end{proof}

Rational data can yield an irrational root. Take two parallel edges
between $s$ and $t$, of resistances $1$ and $a>0$, with unit through-flow.
Orient the first from $s$ to $t$ and the second from $t$ to $s$, and
choose the first as chord. The flows are $q$ and $q-1$;
$\Phi(0)=-a<0<1=\Phi(1)$, and on $(0,1)$ the equation is
$q^2-a(1-q)^2=0$. Its root is
\[
 q=\frac{\sqrt a}{1+\sqrt a},\qquad
 1-q=\frac{1}{1+\sqrt a},\qquad q^2=a(1-q)^2.
\]
For rational $a$ that is not a rational square, $q$ is irrational since
$\sqrt a=q/(1-q)$. Subdividing the second edge into two edges of
resistance $a/2$ gives the same example on a simple triangle.

\begin{lemma}[Degree count]
\label{lem:a-pre-degree}
For a connected block $H$ write $n_H=|V(H)|$, $m_H=|E(H)|$, and
$r_H=m_H-n_H+1$. Then
\[
 \sum_{v\in V(H)}(\deg_H v-2)=2r_H-2.
\]
If $H$ is not a bridge, at most $2r_H-2$ vertices have degree at least three.
\end{lemma}
\begin{proof}
The degree sum is $2m_H$, so subtracting $2n_H$ gives $2r_H-2$.
Every vertex of a nonbridge block has degree at least two, and every
vertex of degree at least three contributes at least one to that sum.
For a bridge, the two degree-one vertices contribute $-2$ in total,
and the second assertion fails at rank zero.
\end{proof}

\subsection{Graph parameters}
\label{subsec:a-pre-parameters}
A cactus is a graph in which distinct simple cycles share at most one
vertex. For multigraphs a pair of parallel edges is a cycle of length two.

\begin{lemma}[Cactus blocks]
\label{lem:a-pre-cactus-blocks}
A connected graph is a cactus if and only if every block is a bridge
or a simple cycle.
\end{lemma}
\begin{proof}
A nonbridge block has minimum degree at least two, so following edges
until a vertex repeats produces a simple cycle $C$. If the block is
not exactly $C$, it has an ear relative to $C$: a path with distinct
endpoints on $C$, with no internal vertex on $C$, and with edges outside
$C$. Indeed, an extra edge between vertices of $C$ is already an
ear. Otherwise a component outside $V(C)$ must have at least two
distinct neighbors on $C$, since a unique neighbor would be a cut
vertex of the block. A path through that component between two such
neighbors is an ear. The ear together with either of the two paths
on $C$ between its endpoints gives two distinct simple cycles sharing
both endpoints. Hence a cactus has no such block.

Conversely, two distinct simple cycles sharing at least two vertices
have a biconnected union: after any one vertex is deleted, each cycle
remains connected and at least one common vertex survives. The union
therefore lies in one block. That block is neither a bridge nor a
simple cycle, since a simple cycle contains no distinct simple cycle.
Thus the stated block condition implies the cactus property.
\end{proof}

Series-parallel means $K_4$-minor-free; a prescribed terminal pair
need not admit a two-terminal series-parallel decomposition. Universal
characterizations by graph class state when the graph must be simple.
The cycle rank of a block $H$ is $r_H=m_H-n_H+1$. Define
\[
 \blockrank=\max_H r_H,\qquad \totalrank=m-n+1,
\]
where the maximum is zero if there are no blocks. For a graph with $k$
connected components the total cycle rank is $m-n+k$.

\begin{lemma}[Total and block cycle rank]
\label{lem:a-pre-ranks}
For a connected graph, $\totalrank=\sum_H r_H$. Consequently bounded
total cycle rank implies bounded block rank, but the converse fails.
Trees have block rank zero. Cacti have block rank at most one, and
exactly one when they contain a cycle. Series-parallel graphs can have
unbounded block rank.
\end{lemma}
\begin{proof}
By \cref{lem:a-pre-incidence}, if there are $h$ blocks and $n>1$, the
block--vertex incidence tree has $n+h$ nodes and $\sum_H n_H$ edges,
so $\sum_H(n_H-1)=n-1$.
Every edge is in exactly one block, so
\[
 \sum_H(m_H-n_H+1)=m-(n-1)=\totalrank.
\]
The one-vertex case has both sides zero. Nonnegativity of block ranks
gives $\blockrank\le\totalrank$. A chain of $k$ triangles joined at
successive cut vertices has block rank one and total rank $k$.
The tree statement follows from its bridge blocks, and the cactus
statement follows from \cref{lem:a-pre-cactus-blocks}. Finally,
$K_{2,k}$, for $k\ge2$, has no $K_4$ minor. Such a minor would require
four disjoint, connected, pairwise adjacent branch sets. At most two
could contain a hub vertex. A connected hub-free branch set must be
a singleton, because the non-hub vertices are independent. The two
hub-free branch sets would therefore be nonadjacent, a contradiction.
Deleting any one vertex of $K_{2,k}$ leaves it connected, so it is one
block; it has $k+2$ vertices, $2k$ edges, and rank $k-1$.
Subdivision preserves this example's rank.
\end{proof}

\subsection{Objectives, scenario domains, and problems}
\label{subsec:a-pre-objectives}
Algorithmic inputs have rational data. A balanced nomination box is the
nonempty compact set
\[
 \nomset=\{b\in\R^V:\ell\le b\le u,\ \one^\top b=0\},
 \qquad \ell,u\in\Q^V,\quad \ell\le u.
\]
Its nonemptiness is equivalent to $\sum_v\ell_v\le0\le\sum_vu_v$:
necessity follows by summation; sufficiency follows by interpolating
between $\ell$ and $u$. Alternatively a nomination family is
$\{b^0+Mz:z\in P\}$, where $b^0,M$ are rational and $P\subseteq\R^k$
is a nonempty compact polytope given by an explicit list of rational
linear equalities and inequalities. We require
$\one^\top(b^0+Mz)=0$ for every $z\in P$. Fixed dimension means that
$k$ is fixed independently of the input length, not that $|V|$ is fixed.
Either nonemptiness is promised, or the algorithm first reports an
empty domain.

Unless specified otherwise, the computational problems use symmetric
quadratic laws. Independent resistance intervals form
\[
 \resset=\prod_{e\in E}[L_e,U_e],\qquad
 L_e,U_e\in\Q,\quad 0<L_e\le U_e<\infty.
\]
Finite resistance uncertainty instead uses
$\resset=\prod_e R_e$ with each $R_e$ an explicitly listed nonempty
finite subset of $\Q_{>0}$. Fixed resistances are singleton domains.
A scenario is $\theta=(b,\beta)\in\Theta:=\nomset\times\resset$ (or the
stated affine-family replacement). All nomination and resistance choices
are independent apart from nomination balance and the stated parameter
polytope. Finite sets are never implicitly replaced by interval hulls.
Asymmetric models specify domains for both positive coefficients.
Other law representations and their encodings are stated where used;
an arbitrary continuous law satisfying the existence theorem is not
thereby a finite algorithmic input.

Write $(x(\theta),\pi(\theta))$ for the unique normalized passive state.
The objectives considered here are
\[
 \begin{aligned}
 F_{s,t}(\theta)&=\pi_s(\theta)-\pi_t(\theta),
 &F_c(\theta)&=c^\top\pi(\theta),\\
 X_e(\theta)&=x_e(\theta), &F_d(\theta)&=d^\top x(\theta).
 \end{aligned}
\]
The ordered pair $(s,t)$ is prescribed. The rational potential weights
satisfy $\one^\top c=0$, ensuring gauge invariance, and their support size
is $p=|\supp c|$. Flow weights $d\in\Q^E$ have no balance requirement.

\begin{lemma}[Continuity and attainment]
\label{lem:a-pre-attainment}
On the scenario domains $\Theta=\nomset\times\resset$ defined above,
including the affine nomination families and positive asymmetric
quadratic coefficients, the map
$\theta\mapsto(x(\theta),\pi(\theta))$ with normalized potentials is
continuous. The same holds on the nomination domains with fixed general
laws, omitting the resistance coordinate. Each of $F_{s,t},F_c,X_e,F_d$
attains its maximum and minimum on its domain.
\end{lemma}
\begin{proof}
For quadratic laws let $\energy_\beta$ denote the primitive energy with
coefficient vector $\beta$: it is $\sum_e\beta_e|x_e|^3/3$ in the
symmetric case and uses $\beta_e^+$ or $\beta_e^-$ according to the sign
of $x_e$ in the asymmetric case. For a convergent sequence
$(b_j,\beta_j)\to(b,\beta)$ with all limiting resistances positive,
spanning-tree routings are uniformly bounded and provide uniformly
bounded comparison energies. A common positive lower bound on the
coefficients makes the minimizing flows uniformly bounded. If a
subsequence converges to $x_*$, then $Ax_*=b$. For any $y$ with $Ay=b$,
add the tree routing of $b_j-b$ to obtain $y_j\to y$ with $Ay_j=b_j$.
Passing to the limit in $\energy_{\beta_j}(x_j)\le
\energy_{\beta_j}(y_j)$ proves that $x_*$ is the unique minimizer at
$(b,\beta)$. Every convergent subsequence therefore has the same limit,
so the flow is continuous. Summing continuous edge drops along a
fixed spanning tree gives continuity of normalized potentials. The
same argument works for fixed general laws, using the coercivity bound
in \cref{thm:a-pre-existence}, and for positive asymmetric coefficients.
The four objectives are continuous functions of the state, and
compactness of each nonempty scenario domain gives attainment of both
maximum and minimum, including for finite resistance sets.
\end{proof}

\begin{definition}[Optimization, validation, and approximation]
\label{def:a-pre-problems}
For any one of the objectives $F$ above, set
$\OPT_F(\Theta)=\max_{\theta\in\Theta}F(\theta)$, which exists by
\cref{lem:a-pre-attainment}.
\begin{enumerate}
\item $\MPD$ is the computation of $\OPT_{F_{s,t}}$ over $b\in\nomset$
with fixed resistances, or jointly over
$(b,\beta)\in\nomset\times\resset$ when stated. Weighted potential
optimization computes $\OPT_{F_c}$; linear flow optimization computes
$\OPT_{F_d}$. No maximization over the terminals is implicit.
\item Arc extremum computation returns
$\min_{\theta\in\Theta}x_e(\theta)$ and
$\max_{\theta\in\Theta}x_e(\theta)$ for a prescribed arc $e$.
Minima may equivalently be computed by maximizing $-X_e$.
\item Given tested arcs $T\subseteq E$ and rational signed capacities
$\underline x_e\le\overline x_e$, robust arc-capacity validation asks
whether
\[
 \forall\theta\in\Theta\ \forall e\in T:\quad
 \underline x_e\le x_e(\theta)\le\overline x_e.
\]
Each scenario's passive state is evaluated before testing the bounds;
capacities do not filter $\Theta$. Its complement is
$\exists\theta\in\Theta\ \exists e\in T$ with
$x_e(\theta)<\underline x_e$ or $x_e(\theta)>\overline x_e$.
\item The exact threshold version of a maximization problem, for rational
$\gamma$, asks
\[
 \OPT_F(\Theta)\ge\gamma,
 \quad\text{equivalently}\quad
 \exists\theta\in\Theta:\ F(\theta)\ge\gamma.
\]
The weak inequality includes equality; it is distinct from a strict
capacity violation. For a minimum, the threshold convention is applied
to $-F$ if it is expressed as a maximization problem.
\item Given rational $\eps>0$, the additive task returns rational
$\underline{v},\overline{v}$ and a rational scenario $\widehat\theta\in\Theta$ such that
\[
 \underline{v}\le\OPT_F(\Theta)\le\overline{v},\qquad
 \overline{v}-\underline{v}\le\eps,\qquad
 F(\widehat\theta)\ge\OPT_F(\Theta)-\eps.
\]
These are certified guarantees, not floating-point estimates. For an
affine family a rational parameter $\widehat z\in P$ may encode the
nomination. Rational points are dense in each of the rational polytopes
in the domain: express a point as a convex combination of rational
vertices and approximate its weights by rational simplex weights.
With continuity from \cref{lem:a-pre-attainment}, applied separately
for each finite resistance choice, this proves that rational
$\eps$-optimal scenarios exist; it gives no polynomial bound on finding
them.
\end{enumerate}
\end{definition}
An additive interval certifies a threshold answer when
$\underline{v}\ge\gamma$ or $\overline{v}<\gamma$.
An interval straddling $\gamma$ need not decide the answer; in
particular, arbitrarily narrow intervals alone do not resolve equality.

\subsection{Computational conventions}
\label{subsec:a-pre-computation}
We use the Turing bit model with total input length $N$: rational data
are binary encoded, and running time includes bit costs and output
length. Polynomial additive time means polynomial in $N$ and
$\max\{0,\log_2(1/\eps)\}$, with the rational encoding of $\eps$ counted
in $N$. Tolerances larger than one may be replaced by one. A bound of
the form $N^{\varphi(k)}$ for parameter $k$ is XP-type (polynomial for fixed
$k$); an FPT bound is $\varphi(k)N^C$ with a constant exponent $C$ independent
of $k$. The same distinction applies with the accuracy bits included.

\paragraph{Square-Root Sum.}
$\SRS$ takes positive integers $a_1,\ldots,a_k,K$ in binary and asks
whether
\[
 \sum_{i=1}^k\sqrt{a_i}\le K.
\]
If zero radicands are allowed,
remove them; an empty sum and a nonpositive threshold are handled
directly. It is open whether $\SRS$ is in P or even in $\NP$;
\citet[Section 1]{etessami2010-on-the-complexity-of-nash} state these open
questions and the $\PSPACE$ upper bound.
\citet[Corollary 1.5]{allender2009-on-the-complexity-of-numerical}
give the stronger counting-hierarchy bound. They use $\ge$; the weak
$\le$ version follows by complementing and accounting for equality.
For positive integer radicands, equality to an integer
occurs only if every radicand is a perfect square: after grouping
square factors, distinct square-free radicals, including $1$, are
linearly independent over $\Q$
\citep[Section 2, Theorem 4]{blomer1991-computing-sums-of-radicals-in}.
For each square-free integer $f>1$ occurring after this grouping, the
coefficient of $\sqrt f$ is positive. Thus equality is testable by
integer square roots. The iterated-PP definition recalled by
\citet[Section 4, citing Wagner and Tor\'an]{allender2009-on-the-complexity-of-numerical}
also gives closure of the counting hierarchy under complement and
polynomial-time Boolean combinations: a deterministic polynomial-time
computation with queries to a fixed level belongs to the next PP-oracle
level. Complementing the weak $\ge$ test and adjoining the equality
case therefore preserves membership in the hierarchy.
\citet[Sections 1 and 7]{blomer1991-computing-sums-of-radicals-in}
give a polynomial-time Monte Carlo test for whether a sum of real
radicals lies in a specified real algebraic number field, including
zero testing, and explicitly distinguish this from determining its
sign; it is not an efficient general ordering algorithm.

\paragraph{Real algebraic computation.}
The class $\ETR$ consists of decision problems polynomial-time reducible
to the existence of a real solution to a finite Boolean combination of
polynomial equalities and inequalities with rational coefficients.
Clearing positive denominators gives integer coefficients. Allowing only
conjunctions gives the same class: Boolean choices and strict inequalities
can be encoded with additional existential variables. For decision
complexity, \citet[Part III, Proposition 2.3 (restating Part I, Proposition 4.1)]{renegar1992-on-the-computational-complexity-and}
constructs all consistent sign vectors of $\sigma$ integer polynomials of
degree at most $D$ in $k$ variables with $(\sigma D)^{O(k)}$ arithmetic
operations and
$\tau\log\tau\log\log\tau\,(\sigma D)^{O(k)}$ bit operations for coefficient
bit bound $\tau$ (replace small $\tau$ by a fixed constant).
Evaluating the input Boolean formula on these vectors decides
existential feasibility, with its formula-evaluation cost included.

We use the following fixed-dimension consequence of
\citet[Theorem 2.27, Section 2.5.2]{basu2014-algorithms-in-real-algebraic-geometry},
which restates the quantifier-elimination theorem from
\citet[Chapter 14]{basu2006-algorithms-in-real-algebraic-geometry}.
For $\sigma$ polynomials of degree at most $D\ge2$ in $k+\nu$ variables,
partition the quantified variables into $\omega$ nonempty blocks of
sizes $k_1,\ldots,k_\omega$. A prenex formula with these blocks and
$\nu$ free variables has an equivalent quantifier-free formula
computable with
\begin{equation}
 \sigma^{(k_\omega+1)\cdots(k_1+1)(\nu+1)}
 D^{O(k_\omega)\cdots O(k_1)O(\nu)}
 \label{eq:a-pre-qe-cost}
\end{equation}
arithmetic operations in the coefficient ring. Output polynomial
degrees are at most $D^{O(k_\omega)\cdots O(k_1)}$; for integer input
coefficients of bit size at most $\tau$, all intermediate and output
integer bit sizes are at most
$\tau D^{O(k_\omega)\cdots O(k_1)O(\nu)}$.
Zero free dimension is handled by adding one unused free variable, so
we take $\nu\ge1$ in these bounds.
The source gives output as a disjunction of conjunctions of disjunctions
of sign tests, with the outer count bounded by
\cref{eq:a-pre-qe-cost}, the conjunction count by
$\sigma^{(k_\omega+1)\cdots(k_1+1)}D^{O(k_\omega)\cdots O(k_1)}$, and the
inner count by $D^{O(k_\omega)\cdots O(k_1)}$.
Thus for fixed total dimension the arithmetic work and coefficient bit
sizes are polynomial in $\sigma,D,\tau$, and ordinary integer arithmetic
turns these into polynomial bit costs; reading and evaluating the
Boolean input formula adds its encoding length. These bounds use
numerical degree and dense polynomial encoding; they do not give
polynomial time for sparse binary exponents or for an unbounded number
of real variables.

\paragraph{Fixed-dimension sample points.}
\label{prim:a-pre-sample-points}
Given an explicit quantifier-free formula in a fixed number $k$ of real
variables, using integer polynomials of degree at most $D$ and coefficient
bit size at most $\tau$, one can compute in polynomial bit time a finite
set of sample points meeting every semialgebraically connected component
of every realizable sign condition of the input polynomials. Testing
the formula at these points supplies a point in every
nonempty semialgebraic set it defines. Each sample is given by a \emph{real
univariate representation}: a common defining univariate polynomial,
an isolating interval selecting its real root $\theta_{\rm alg}$, and
all coordinates as rational functions of that root, with denominators
nonzero there. Degrees and encoding lengths are polynomial for fixed
$k$. See \citet[Chapter 13, Sections 13.1 (finding realizable sign
conditions) and 13.3 (sample points), and Chapter 12, Section 12.4
(univariate representations)]{basu2006-algorithms-in-real-algebraic-geometry}.
All coordinates of a sample lie in the single real algebraic field
$\Q(\theta_{\rm alg})$, of polynomial degree; this is what a
\emph{common algebraic representation} means here. The subscript
separates the algebraic generator from a scenario $\theta$.

\paragraph{Exact and additive outputs.}
A real algebraic number $\alpha$ is represented by a nonzero square-free
integer polynomial $P_\alpha$ and rational endpoints $\eta_-<\eta_+$
with exactly one real root of $P_\alpha$ in $(\eta_-,\eta_+)$, namely
$\alpha$, and no endpoint root. A Thom encoding (signs of successive
derivatives at the root) is an alternative
\citep[Chapter 10]{basu2006-algorithms-in-real-algebraic-geometry}. The size counts the polynomial's degree, all coefficient
bits, and the endpoint or encoding bits. For a degree-$D$ dense
polynomial with coefficient bit bound $\tau$ and endpoints of total
bit length $\kappa$, it is $O((D+1)(\tau+1)+\kappa)$.
The polynomial need not be minimal. Repeated roots can be removed using
$\gcd(P_\alpha,P_\alpha')$. Refining an isolating interval by rational bisection and
exact univariate root counts permits rational rounding to prescribed
accuracy, with bit costs measured in the representation size and
accuracy bits. Root counting via Sturm sequences and sign determination
are described in
\citet[Chapters 2 and 10]{basu2006-algorithms-in-real-algebraic-geometry};
see also \citet[Section 2.5.1]{basu2014-algorithms-in-real-algebraic-geometry}
for sign determination.
Algebraic output claims apply to semialgebraic input models; arbitrary
continuous laws need not give algebraic states or values.

The output contract distinguishes three objects.
\begin{itemize}
\item An exact value is the optimum itself in the stated exact encoding;
an additive value is a certified rational interval containing it.
A sum of separately encoded block values is an expression encoding,
not automatically a single polynomial-size isolating representation
or a polynomial-time exact comparison procedure.
\item An exact optimizing scenario is an admissible parameter choice
attaining the optimum, represented by rational or algebraic coordinates
as specified. An additive optimizing scenario is an exactly admissible
rational choice with the objective guarantee in
\cref{def:a-pre-problems}; it need not be close in coordinates to an
exact optimizer.
\item An exact state satisfies both passive equations exactly, with
normalized potentials. An additive state output consists of rational
coordinate intervals of width at most $\eps$ containing that scenario's
unique exact state, or rational coordinates with certified sup-norm
error at most $\eps$. Rounded coordinates need not themselves satisfy
the equations; a small residual alone is not this error guarantee.
\end{itemize}
Rational scenarios may have irrational states and values, as the
parallel example shows. No output kind is inferred from another without
an explicit argument and a bound on its encoding size.

\subsection{Notation}
\label{subsec:a-pre-notation}
\begin{center}
\small
\begin{tabular}{@{}ll@{}}
\toprule
Symbol & Meaning \\
\midrule
$G=(V,E),\ n,m$ & Oriented network and its vertex and edge counts \\
$A,\one$ & Tail-positive incidence matrix; all-ones vector \\
$b,x,\pi$ & Balanced nominations, signed flows, potentials \\
$g_e,\beta_e,\beta_e^\pm$ & Edge law and positive quadratic coefficients \\
$\energy,\ b^\top\pi$ & Primitive energy; constitutive dissipation \\
$H,b^H,r_H$ & Block, aggregated nomination, cycle rank \\
$\varkappa,\lambda,q$ & Incident-block count, spanning-path flow, circulation \\
$\blockrank,\totalrank$ & Maximum block rank; total cycle rank \\
$\nomset,\resset,\Theta$ & Nomination, resistance, and scenario domains \\
$(s,t),c,p,d$ & Ordered terminals, potential weights/support, flow weights \\
$b^0,M,z,P$ & Affine offset, matrix, parameter, parameter polytope \\
$\ell,u,L_e,U_e$ & Nomination bounds; resistance bounds \\
$\OPT_F,\gamma,\eps$ & Maximum objective value, threshold, additive tolerance \\
$\underline{v},\overline{v}$ & Certified additive value interval endpoints \\
$\sigma,D,\tau,\nu$ & Polynomial count, degree/bit bounds, free-variable count \\
$f,\kappa$ & Square-free integer; isolating-endpoint bit length \\
$N,\SRS,\ETR$ & Input bit length, Square-Root Sum, existential real class \\
\bottomrule
\end{tabular}
\end{center}

The distinction between block rank and total cycle rank is central
below. For a prescribed terminal pair, the block path decomposes the
objective into local drops, which must still be combined in the Turing
model. The next two sections use nomination structure and separate
approximation of local values to avoid a potentially large algebraic
field. Later sections examine which parts of this approach survive
changes in edge laws, resistance domains, and objectives.
